\documentclass[11pt,a4paper]{article}
\usepackage[utf8]{inputenc}
\usepackage[T1]{fontenc}
\usepackage[ngerman]{babel}
\usepackage{helvet}
\renewcommand{\familydefault}{\sfdefault}
\usepackage[left=2cm, right=2cm, top=2cm, bottom=2cm]{geometry}
\usepackage{amsmath,amssymb}
\usepackage{array}
\usepackage{booktabs}
\usepackage{tikz}
\usepackage{pgfplots}
\pgfplotsset{compat=1.18}
\usepackage{fancyhdr}
\usepackage{xcolor}

% Lösungsfarbe
\definecolor{correct}{RGB}{34,139,94}

\pagestyle{fancy}
\fancyhf{}
\fancyhead[L]{\small Physik Klasse 10E}
\fancyhead[R]{\small\bfseries ERWARTUNGSHORIZONT -- Variante C}
\fancyfoot[C]{\small Seite \thepage}
\renewcommand{\headrulewidth}{0.4pt}

\setlength{\parindent}{0pt}
\setlength{\fboxsep}{8pt}

% Aufgaben-Befehl im Layout B Stil
\newcounter{aufgabe}
\newcommand{\aufgabe}[2]{%
\stepcounter{aufgabe}%
\vspace{0.5cm}%
\noindent\fbox{\parbox{\dimexpr\textwidth-2\fboxsep-2\fboxrule}{%
\textbf{Aufgabe \theaufgabe: #1}\hfill\textbf{#2 BE}%
}}%
\vspace{0.3cm}%
}

\newcommand{\lsg}[1]{\textcolor{correct}{#1}}
\newcommand{\be}[1]{\hfill\textcolor{correct}{\small(#1 BE)}}

\begin{document}

% === KOPFBEREICH ===
\noindent\fbox{\parbox{\dimexpr\textwidth-2\fboxsep-2\fboxrule}{
\centering
{\Large\bfseries Probeklausur: Kinematik -- Variante C}\\[0.2cm]
{\large\bfseries Erwartungshorizont}\\[0.2cm]
{\small Klasse 10E \quad|\quad 90 Minuten \quad|\quad 60 BE}
}}

\vspace{0.4cm}

% === NOTENSPIEGEL ===
\begin{center}
\small
\begin{tabular}{|c|c|c|c|c|c|c|c|c|c|c|c|c|c|c|c|}
\hline
\textbf{NP} & 15 & 14 & 13 & 12 & 11 & 10 & 9 & 8 & 7 & 6 & 5 & 4 & 3 & 2 & 1 \\
\hline
\textbf{ab BE} & 60 & 59 & 58 & 55 & 51 & 48 & 44 & 41 & 36 & 32 & 29 & 24 & 20 & 17 & 12 \\
\hline
\end{tabular}

\vspace{0.2cm}
{\footnotesize NP = Notenpunkte (15-Punkte-Skala, Sek II MV)}
\end{center}

\vspace{0.3cm}

% === LÖSUNGEN ===

\aufgabe{Grundlagen}{6}

\begin{enumerate}
\item[a)] \lsg{36 km/h = 36 $\div$ 3,6 = 10 m/s} \be{1}

\item[b)] \lsg{Formel: $s = \frac{1}{2}gt^2$ $\Rightarrow$ $t = \sqrt{\frac{2s}{g}}$} \be{1}\\
\lsg{Einsetzen: $t = \sqrt{\frac{2 \cdot 45\,\mathrm{m}}{10\,\mathrm{m/s^2}}} = \sqrt{9\,\mathrm{s^2}} = 3\,\mathrm{s}$} \be{1}

\item[c)] Tabelle: \be{3}

\begin{center}
\begin{tabular}{|l|c|c|}
\hline
\textbf{Diagramm} & \textbf{Steigung bedeutet} & \textbf{Fläche bedeutet} \\
\hline
$s(t)$-Diagramm & \lsg{Geschwindigkeit $v$} & --- \\
\hline
$v(t)$-Diagramm & \lsg{Beschleunigung $a$} & \lsg{Weg $s$} \\
\hline
\end{tabular}
\end{center}
{\small\color{correct} Je 1 BE pro richtige Eintragung}
\end{enumerate}

\aufgabe{Motorrad-Bremsung}{8}

\begin{enumerate}
\item[a)] \lsg{$v_0 = 72\,\mathrm{km/h} = 72 \div 3{,}6 = 20\,\mathrm{m/s}$} \be{1}

\item[b)] \lsg{Formel: $v = v_0 + a \cdot t$ $\Rightarrow$ $t = \dfrac{v_0}{|a|}$} \be{1}\\[0.2cm]
\lsg{Einsetzen: $t = \dfrac{20\,\mathrm{m/s}}{4\,\mathrm{m/s^2}} = 5\,\mathrm{s}$} \be{1}

\item[c)] \lsg{Formel: $s = \dfrac{v_0^2}{2|a|}$} \be{1}\\[0.2cm]
\lsg{Einsetzen: $s = \dfrac{(20\,\mathrm{m/s})^2}{2 \cdot 4\,\mathrm{m/s^2}} = \dfrac{400}{8}\,\mathrm{m} = 50\,\mathrm{m}$} \be{1}

\item[d)] \lsg{Nein, der Fahrer kann nicht rechtzeitig anhalten.} \be{1}\\
\lsg{Begründung: Der Bremsweg (50 m) ist länger als der Abstand zum Hindernis (40 m).} \be{1}\\
\lsg{Es fehlen 10 m: $50\,\mathrm{m} - 40\,\mathrm{m} = 10\,\mathrm{m}$} \be{1}
\end{enumerate}

\aufgabe{Volleyball-Aufschlag}{8}

\begin{enumerate}
\item[a)] \lsg{Formel: $v = v_0 - g \cdot t$ $\Rightarrow$ $t_{\text{steig}} = \dfrac{v_0}{g}$} \be{1}\\[0.2cm]
\lsg{Einsetzen: $t_{\text{steig}} = \dfrac{10\,\mathrm{m/s}}{10\,\mathrm{m/s^2}} = 1\,\mathrm{s}$} \be{1}

\item[b)] \lsg{Formel: $h_{\text{max}} = \dfrac{v_0^2}{2g}$} \be{1}\\[0.2cm]
\lsg{Einsetzen: $h_{\text{max}} = \dfrac{(10\,\mathrm{m/s})^2}{2 \cdot 10\,\mathrm{m/s^2}} = \dfrac{100}{20}\,\mathrm{m} = 5\,\mathrm{m}$} \be{2}

\item[c)] \lsg{Der Ball kommt mit $v = 10\,\mathrm{m/s}$ (nach unten) an.} \be{1}\\
\lsg{Begründung: Energieerhaltung / Symmetrie der Bewegung / zeitlose Formel liefert $|v| = v_0$} \be{2}
\end{enumerate}

\newpage

\aufgabe{Autofahrt in der Stadt}{8}

\begin{enumerate}
\item[a)] \lsg{Phase A (0--4 s): Beschleunigte Bewegung, das Auto beschleunigt von 0 auf 12 m/s.} \be{1}\\
\lsg{Phase B (4--16 s): Gleichförmige Bewegung mit konstant 12 m/s.} \be{1}\\
\lsg{Phase C (16--20 s): Bremsvorgang, das Auto bremst von 12 m/s auf 0 ab.} \be{1}

\item[b)] \lsg{Formel: $a = \dfrac{\Delta v}{\Delta t}$} \be{1}\\[0.2cm]
\lsg{Einsetzen: $a = \dfrac{12\,\mathrm{m/s} - 0}{4\,\mathrm{s}} = 3\,\mathrm{m/s^2}$} \be{1}

\item[c)] \lsg{Fläche unter dem Graphen = Weg. Zerlegung in geometrische Figuren:} \be{1}\\
\lsg{Phase A (Dreieck): $s_A = \dfrac{1}{2} \cdot 4\,\mathrm{s} \cdot 12\,\mathrm{m/s} = 24\,\mathrm{m}$}\\
\lsg{Phase B (Rechteck): $s_B = 12\,\mathrm{s} \cdot 12\,\mathrm{m/s} = 144\,\mathrm{m}$}\\
\lsg{Phase C (Dreieck): $s_C = \dfrac{1}{2} \cdot 4\,\mathrm{s} \cdot 12\,\mathrm{m/s} = 24\,\mathrm{m}$}\\
\lsg{Gesamtweg: $s = 24 + 144 + 24 = 192\,\mathrm{m}$} \be{2}
\end{enumerate}

\aufgabe{Fallexperimente}{8}

\begin{enumerate}
\item[a)] \lsg{Im Vakuum gibt es keinen Luftwiderstand.} \be{1}\\
\lsg{Beide Körper werden nur durch die Schwerkraft beschleunigt.} \be{1}\\
\lsg{Die Fallbeschleunigung $g$ ist für alle Körper gleich, unabhängig von der Masse.} \be{1}

\item[b)] \lsg{In Luft wirkt Luftwiderstand: $F_W \sim v^2 \cdot A$.} \be{1}\\
\lsg{Die Feder hat eine viel größere Fläche im Verhältnis zur Masse.} \be{1}\\
\lsg{Daher erreicht die Feder schneller ihre (niedrigere) Grenzgeschwindigkeit.} \be{1}

\item[c)] \lsg{Formel: $v = \sqrt{2gs}$}\\
\lsg{Einsetzen: $v = \sqrt{2 \cdot 10\,\mathrm{m/s^2} \cdot 2000\,\mathrm{m}} = \sqrt{40000\,\mathrm{m^2/s^2}} = 200\,\mathrm{m/s}$} \be{1}\\
\lsg{Tatsächlich viel geringer, da der Regentropfen seine Grenzgeschwindigkeit erreicht.} \be{1}
\end{enumerate}

\aufgabe{Weg-Zeit-Gesetz aufstellen}{10}

\begin{enumerate}
\item[a)] \lsg{Formel: $s(t) = v \cdot t + s_0$}\\
\lsg{Einsetzen: $s(t) = 25\,\mathrm{m/s} \cdot t + 50\,\mathrm{m}$} \be{2}

\item[b)] \lsg{Einsetzen: $s(8) = 25 \cdot 8 + 50 = 200 + 50 = 250\,\mathrm{m}$} \be{2}

\item[c)] \lsg{$300 = 25t + 50$ $\Rightarrow$ $250 = 25t$ $\Rightarrow$ $t = 10\,\mathrm{s}$} \be{2}

\item[d)] \lsg{Gerade von $(0; 50)$ nach $(10; 300)$:} \be{2}

\begin{center}
\begin{tikzpicture}
\begin{axis}[
    width=8cm, height=4cm,
    xlabel={$t$ in s},
    ylabel={$s$ in m},
    xmin=0, xmax=10,
    ymin=0, ymax=350,
    grid=major,
    axis lines=left,
    xtick={0,2,4,6,8,10},
    ytick={0,50,100,150,200,250,300,350}
]
\addplot[thick, correct] coordinates {(0,50) (10,300)};
\end{axis}
\end{tikzpicture}
\end{center}

\item[e)] \lsg{Zug 1: $s_1(t) = 25t + 50$, Zug 2: $s_2(t) = t^2$}\\
\lsg{Gleichsetzen: $t^2 = 25t + 50$ $\Rightarrow$ $t^2 - 25t - 50 = 0$}\\
\lsg{$t = \dfrac{25 + \sqrt{825}}{2} \approx 26{,}9\,\mathrm{s}$} \be{2}
\end{enumerate}

\newpage

\aufgabe{Freier Fall mit Anfangsgeschwindigkeit}{10}

\begin{enumerate}
\item[a)] \lsg{Formel: $s = v_0 \cdot t + \frac{1}{2}gt^2$} \be{1}\\
\lsg{Einsetzen: $80 = 10t + 5t^2$ $\Rightarrow$ $t^2 + 2t - 16 = 0$}\\
\lsg{$t = \dfrac{-2 + \sqrt{68}}{2} \approx 3{,}1\,\mathrm{s}$} \be{2}

\item[b)] \lsg{Formel: $v = v_0 + g \cdot t$} \be{1}\\
\lsg{Einsetzen: $v = 10 + 10 \cdot 3{,}1 = 41{,}1\,\mathrm{m/s}$} \be{1}

\item[c)] \lsg{Freier Fall: $80 = 5t^2$ $\Rightarrow$ $t = 4\,\mathrm{s}$} \be{2}\\
\lsg{Zeitersparnis: $\Delta t = 4{,}0 - 3{,}1 = 0{,}9\,\mathrm{s}$} \be{1}

\item[d)] \lsg{Die Erdbeschleunigung erhöht die Geschwindigkeit um 10 m/s pro Sekunde.} \be{1}\\
\lsg{Nach 3 s ist der Ball bereits ca. 40 m/s schnell -- die anfänglichen 10 m/s fallen kaum ins Gewicht.} \be{1}
\end{enumerate}

\vspace{0.3cm}
\hrule
\vspace{0.3cm}

\noindent\fbox{\parbox{\dimexpr\textwidth-2\fboxsep-2\fboxrule}{%
\textbf{Zusatzaufgabe} (max. 2 Zusatz-BE)
}}

\vspace{0.3cm}

\lsg{Bahn A: $s_A(t) = 15\,\mathrm{m/s} \cdot t$}\\
\lsg{Bahn B (100 m dahinter): $s_B(t) = t^2 - 100\,\mathrm{m}$}\\[0.2cm]
\lsg{Gleichsetzen: $t^2 - 100 = 15t$ $\Rightarrow$ $t^2 - 15t - 100 = 0$}\\
\lsg{$(t - 20)(t + 5) = 0$ $\Rightarrow$ $t = 20\,\mathrm{s}$} \be{2}

{\small\color{correct} Probe: $s_A(20) = 300\,\mathrm{m}$, $s_B(20) = 300\,\mathrm{m}$ \checkmark}

\vspace{0.3cm}
\hrule
\vspace{0.3cm}

\begin{center}
\fbox{\parbox{0.8\textwidth}{
\centering
\textbf{Zusammenfassung der Anforderungsbereiche}\\[0.3cm]
\begin{tabular}{|l|c|c|c|}
\hline
\textbf{AFB} & \textbf{Aufgaben} & \textbf{BE} & \textbf{Anteil} \\
\hline
I (Reproduktion) & 1 & 6 & 10\% \\
\hline
II (Anwendung) & 2, 3, 4, 6a--d & 32 & 53\% \\
\hline
III (Transfer) & 5, 6e, 7, Zusatz & 20 (+2) & 37\% \\
\hline
\textbf{Gesamt} & & \textbf{58 (+2)} & 100\% \\
\hline
\end{tabular}
}}
\end{center}

\end{document}
